\chapter{جمع‌بندی و کارهای آینده}
\label{ch:conclusion}

\section{خلاصه‌ی کار انجام‌شده}
\label{sec:conclusion-summary}

در این پروژه یک مقیاس‌گذار خودکار افقی به زبان \lr{Go} طراحی و پیاده‌سازی شد که سیگنال
مقیاس‌گذاری را با یک عبارت \lr{PromQL} پیکربندی‌پذیر از \lr{Prometheus} می‌خواند، اندازه‌ی
ناوگان را با یکی از سه سیاست تصمیم می‌گیرد، و آن را از راه زیرمنبع \lr{scale} روی
\lr{Kubernetes} اعمال می‌کند. معماری در امتداد سه دغدغه‌ی مستقل شکسته شد و هر مرحله واسطی با
دست‌کم دو پیاده‌سازی دارد.

سپس این مقیاس‌گذار روی یک خوشه‌ی واقعی و در برابر چهار الگوی بار و نُه بازوی مقایسه ارزیابی
شد. ارزیابی شامل ۴۹ اجرای معتبر با مقیاس زمانی حقیقی است که هر کدام ۱۸۰۰ ثانیه اندازه‌گیری
دارند، به‌علاوه‌ی یک چارچوب آزمایش با بررسی اعتبار خودکار که اجراهای معیوب را پیش از ورود به
تحلیل کنار می‌گذارد.

اما همان‌گونه که در بخش~\ref{sec:intro-question} گفته شد، خود مقیاس‌گذار \emph{ابزار} این
پژوهش است و نه دستاورد آن. دستاورد، پاسخ‌هایی است که در ادامه می‌آید.

\section{پاسخ به پرسش‌های پژوهش}
\label{sec:conclusion-answers}

\subsection*{پرسش ۱: کدام سیگنال حلقه را ناپایدار می‌کند؟}

پاسخ ساختاری است و نه صرفاً تجربی. پیش‌بینی باید روی سیگنالی انجام شود که
\textbf{برون‌زا} باشد؛ یعنی سیگنالی که کنش‌های خود کنترل‌گر آن را جابه‌جا نکند.

بار هر نمونه، $m(t) = \lambda(t)/r(t)$، درون‌زاست چون $r(t)$ خروجی خود کنترل‌گر است.
برون‌یابی آن یک مسیر بسته با بهره‌ی مثبت می‌سازد: مقیاس به بالا باعث افت بار هر نمونه
می‌شود، پیش‌بینی‌کننده این افت را ادامه‌دار می‌بیند، مقیاس به پایین توصیه می‌کند، و چرخه از
نو آغاز می‌شود.

نرخ کل ورود درخواست‌ها، $\lambda(t)$، برون‌زاست چون مقیاس‌گذار نمی‌تواند تعداد درخواست‌های
کاربران را تغییر دهد. سازمایه‌ی این تفاوت را می‌توان مستقیماً در فرمول دید: در
$d = \lceil \hat{\lambda}/T \rceil$ هیچ وابستگی‌ای به $r$ وجود ندارد، پس بهره‌ی آن مسیر
\emph{ساختاراً} صفر است.

این پاسخ در سه سطح مستقل تأیید شد: در آزمون واحد حلقه‌بسته (۲۹ تغییر در برابر صفر، زیر بار کل
ثابت)، در شبیه‌ساز (۱۷ تغییر در برابر ۳، با ۲۳ درصد هزینه‌ی بیشتر و بدون سود)، و در خوشه‌ی
واقعی (بخش~\ref{sec:results-stabilizer}).

مشاهده‌ی تکمیلی و غیرمنتظره‌ی این پروژه آن است که تحریک‌کننده‌ی ناپایداری، \emph{ناپیوستگی در
ورودی یا در مشتق آن} است و نه صرفاً برگشتن بار. تنها الگوی هموارِ مجموعه، تنها الگویی است که
آسیب‌شناسی در آن ظاهر نمی‌شود.

\subsection*{پرسش ۲: میراگر چه مقدار را می‌پوشاند؟}

تقریباً همه‌اش را، و این نتیجه‌ی اصلی این پایان‌نامه است.

با پنجره‌ی پایدارسازی ۹۰ ثانیه‌ای، سیاست معیوب روی \arm{ramp} صفر درصد و روی \arm{spike} ۰/۶
درصد نقض کیفیت خدمت دارد و روی \arm{diurnal} حتی از سیاست درست \emph{ارزان‌تر} است. با
خاموش کردن همان میراگر و بدون هیچ تغییر دیگری، همان سیاست به ۳۹ تا ۵۱ درصد نقض و ۹۸ تا ۱۱۹
تغییر جهت در ۱۲۰ چرخه‌ی کنترل می‌رسد و ناوگان خود را تا یک نمونه فرو می‌پاشاند.

\begin{quote}
\textbf{یک پنجره‌ی پایدارسازی ۹۰ ثانیه‌ای می‌تواند یک قانون کنترل بنیاداً معیوب را چنان
بپوشاند که از همه‌ی آزمون‌های کیفیت خدمت و هزینه سربلند بیرون بیاید.}
\end{quote}

بهایی که این پوشاندن دارد در سنجه‌های متعارف پیدا نیست، و همین است که آن را خطرناک می‌کند.
پیامد عملی برای هر کسی که یک سیاست مقیاس‌گذاری را ارزیابی می‌کند این است: \textbf{ارزیابی
پایداری باید میراگر را خاموش کند و بار کاری با ناپیوستگی بدهد.}

\subsection*{پرسش ۳: پیش‌بینی چه می‌خرد، و بهایش چیست؟}

پاسخ یک \textbf{مبادله} است و نه یک برتری.

پیش‌بینی، کیفیت خدمت و سرعت واکنش می‌خرد: در میان بازوهای خودکار، کم‌ترین کم‌تأمین منابع
روی هر چهار بار کاری و کم‌ترین یا هم‌تراز کم‌ترین نقض، و زمان واکنش ۲۰ تا ۳۵ ثانیه در برابر
۳۲/۵ تا ۱۲۷/۵ ثانیه برای \arm{hpa-custom} که همان سیگنال و همان نقطه‌ی هدف را دارد. بهای
آن، ۵/۷ تا ۲۰/۲ درصد ظرفیت بیشتر است.

این پاسخ در جریان کار یک بار \emph{قوی‌تر} بیان شده بود و پس گرفته شد. آن بیان قوی‌تر بر
داده‌ای استوار بود که یک نقص در مولد بار تولید کرده بود؛ اندازه‌گیری مجدد بازوی مقایسه نشان
داد که برتری در محور هزینه وجود ندارد (بخش~\ref{sec:results-hpacustom-redo}).

پرسش سوم یک بخش دیگر هم داشت: سود پیش‌بینی روی کدام جنس بار کاری وجود دارد؟ پاسخ، یک
\textbf{نتیجه‌ی منفی نسبت به پیشنهاد اولیه‌ی پروژه} است. پیشنهاد مدعی بود پیش‌بینی به
«جهش‌های ناگهانی ترافیک» کمک می‌کند. این ادعا نادرست است و اندازه‌گیری خلافش را نشان می‌دهد:
هیچ پیش‌بینی‌کننده‌ای یک پله‌ی آنی را از پیش نمی‌بیند، و روی \arm{spike} سیاست پیش‌بینانه در
همان لحظه‌ای واکنش نشان می‌دهد که سیاست آستانه‌ای.

ادعای قابل دفاع این است: \textbf{پیش‌بینی روی بارهای روندار و دوره‌ای کمک می‌کند و روی
پله‌های پیش‌بینی‌ناپذیر ذاتاً نمی‌تواند کمک کند.} این بیان ضعیف‌تر از ادعای پیشنهاد است و
دقیقاً به همین دلیل قابل دفاع است.

\section{درس‌های مهندسی}
\label{sec:conclusion-lessons}

جدا از پاسخ‌های بالا، چند درس عملی در جریان کار به دست آمد که ارزش ثبت مستقل دارند. وجه
مشترک همه‌ی آن‌ها این است که به‌جای شکست پرسروصدا، نتایج را \emph{بی‌سروصدا} خراب می‌کنند.

\textbf{یک مقیاس‌گذار مبتنی بر سنجه‌های برنامه، دقیقاً وقتی نابینا می‌شود که بیش از همه به آن
نیاز است.} زیر بار اضافی، محدودیت پردازنده کانتینر را منجمد می‌کند، نقطه‌ی سنجه‌ها پاسخ
نمی‌دهد، پرس‌وجو بردار تهی برمی‌گرداند، و مقیاس‌گذار آن را \mdash{} به‌درستی \mdash{} صفر می‌خواند
(بخش~\ref{sec:method-observability-collapse}). هیچ مؤلفه‌ای اشتباه نمی‌کند و مجموع شکست
می‌خورد.

\textbf{چارچوب آزمایش خودش یک دستگاه اندازه‌گیری است.} سه نقص در کد \emph{بررسی‌کننده} پیدا
شد و نه در سامانه‌ی تحت آزمون، و هر سه در جهت رد کردن اجراهای \emph{سالم} سوگیری داشتند
(بخش~\ref{sec:method-instrument}). دستگاهی که هرگز روی ورودی معلوم خوانده نشده، کالیبره نشده
است.

\textbf{یک گیت که به دلیل اشتباه شلیک می‌کند، از نبودِ گیت فقط اندکی بهتر است}، چون نفر بعدی
را به سراغ اشکال‌زدایی مؤلفه‌ی درست نمی‌فرستد.

\textbf{یک برچسب، یک هویت نیست} (بخش~\ref{sec:results-image}). تنها مرجع قابل اتکا برای اینکه
چه چیزی اجرا شده، اثرانگشت ثبت‌شده در وضعیت خود نمونه است.

\textbf{یک بازوی مقایسه‌ی معیوب، بازوی خودی را رایگان جلوه می‌دهد}
(بخش~\ref{sec:results-hpacustom-redo}). این دقیقاً همان حالتی است که آزمون حذفی برای گرفتن آن
وجود دارد، و تنها به این دلیل گرفته شد که آن بازو دوباره اجرا شد و به آن اعتماد نشد.

\textbf{نوشتن آزمون به‌مثابه‌ی مشخصه، نقص‌هایی را آشکار می‌کند که خواندن کد آشکار نمی‌کند}
(بخش~\ref{sec:design-defects}). هر سه نقص مسیر کنترل تنها در بدترین وضعیت خوشه فعال می‌شدند.

\section{محدودیت‌ها}
\label{sec:conclusion-limitations}

محدودیت‌های اعتبار به‌تفصیل در بخش~\ref{sec:results-validity} آمده‌اند و در اینجا فقط فهرست
می‌شوند تا خواننده بتواند دامنه‌ی تعمیم نتایج را بسنجد:

\begin{itemize}
  \item خوشه تک‌گره و بدون رقابت با مستأجران دیگر است.
  \item تأخیر در بازه‌ی بار این آزمایش تفکیک‌کننده نیست؛ ظرفیت واقعی هر نمونه بیش از بیست
  برابر نرخ هدف است.
  \item سنجه‌ی نقض کیفیت خدمت، رزرو ظرفیت را می‌سنجد و نه آسیب مشاهده‌شده‌ی کاربر.
  \item تنها چند بازو سه بار تکرار شده‌اند؛ قدرت تفکیک حدود ۱/۵ واحد درصد است.
  \item تنها یک سرویس بدون‌حالت آزموده شده و هیچ زنجیره‌ی چندسرویسی در کار نیست.
  \item تصویر اندازه‌گیری‌شده از کد نهایی مخزن ساخته‌شدنی نیست؛ تثبیت و مستند شده است.
\end{itemize}

\section{کارهای آینده}
\label{sec:conclusion-future}

\textbf{اشتقاق تحلیلی مرز پایداری.} در این نوشتار، شرط پایداری به‌صورت \emph{ساختاری} بیان شد
(برون‌زا در برابر درون‌زا) و به‌صورت تجربی تأیید شد. گام بعدی، به دست آوردن یک شرط کمّی بر
حسب افق پیش‌بینی $h$، ضریب هموارسازی $\alpha$، فاصله‌ی آشتی‌دهی و تأخیر راه‌اندازی است، و
سپس اعتبارسنجی مرز پیش‌بینی‌شده در برابر یک جاروب پارامتری در شبیه‌ساز. با توجه به آنچه
بخش~\ref{sec:results-simulator} درباره‌ی محدودیت شبیه‌ساز روی بارهای پله‌ای گفت، هر نتیجه‌ی
چنین جاروبی به وارسی نقطه‌ای روی خوشه نیاز دارد.

\textbf{جاروب سود پیش‌بینی در برابر تأخیر راه‌اندازی.} تمام استدلال این پروژه بر نسبت میان
«سرعت حرکت بار» و «زمان آماده‌شدن نمونه» استوار است. تغییر دادن سیستماتیک عامل دوم \mdash{} که با
تزریق تأخیر مصنوعی در کاوشگر آمادگی ساده است \mdash{} یک یافته‌ی تعمیم‌پذیر می‌دهد: از چه تأخیری به
بعد پیش‌بینی می‌ارزد.

\textbf{مدل هزینه و قاعده‌ی تصمیم سربه‌سر.} با توجه به اینکه نتیجه یک مبادله است، پرسش عملی
بعدی این است که با چه نسبتی میان هزینه‌ی ظرفیت و هزینه‌ی نقض کیفیت خدمت، سیاست پیش‌بینانه
انتخاب درستی است. داده‌های موجود برای ساختن چنین مدلی کافی است.

\textbf{افزودن \lr{KEDA} به‌عنوان بازوی پنجم.} این کار مقایسه را به رایج‌ترین ابزار عملی
بوم‌سازگان گسترش می‌دهد.

\textbf{پیش‌بینی‌کننده‌های سنگین‌تر.} \lr{ARIMA} یا \lr{LSTM} پشت همان واسط
\code{Forecaster}. باید تأکید کرد که این کار تنها با رعایت یافته‌ی این پروژه معنا دارد: هر
پیش‌بینی‌کننده‌ی جدید باید نرخ ورود \emph{برون‌زا} را پیش‌بینی کند و نه سیگنالی را که
کنترل‌گر جابه‌جا می‌کند. پیش‌بینی‌کننده‌ی دقیق‌تر روی سیگنال نادرست، تنها ناپایداری را
سریع‌تر می‌کند.

\textbf{اثر رگبار سرد.} پدیده‌ای که در بخش~\ref{sec:results-otherlimits} گزارش شد،
بازتولیدپذیر است اما سازوکارش ناشناخته است. اگر تأیید شود، مسیری است که از راه آن پیش‌بینی
می‌تواند سودی بدهد که سنجه‌های فعلی این پروژه اصلاً نمی‌بینند: نمونه‌ای که پیش از رسیدن بار
افزوده شود گرم است و نمونه‌ای که واکنشی افزوده شود سرد.

\textbf{زنجیره‌های چندسرویسی.} در یک زنجیره، مقیاس‌گیری یک سرویس بار سرویس بعدی را تغییر
می‌دهد. این یک مسیر پس‌خور تازه است، و با توجه به یافته‌ی مرکزی این پروژه، دقیقاً همان جنس
ساختاری است که ارزش بررسی دارد.

\textbf{آشتی دادن تصویر و کد مخزن}، مطابق آنچه در بخش~\ref{sec:results-image} توضیح داده شد.

\section{سخن پایانی}
\label{sec:conclusion-final}

این پروژه با هدف ساختن یک مقیاس‌گذار خودکار سفارشی آغاز شد و با یافته‌ای درباره‌ی
\emph{ارزیابی} مقیاس‌گذارها پایان یافت. تغییر مسیر از یک اندازه‌گیری در هفته‌ی نخست آمد که
نشان داد نمونه‌ی اولیه سیگنال نادرستی را پیش‌بینی می‌کند \mdash{} و مهم‌تر از آن، نشان داد که این
خطا در سنجه‌های متعارف تقریباً نامرئی است.

آنچه در نهایت به‌عنوان مهم‌ترین نتیجه ثبت شد، درباره‌ی هیچ‌یک از سیاست‌ها نیست، بلکه درباره‌ی
روشی است که با آن سیاست‌ها سنجیده می‌شوند: سازوکاری که برای بهبود رفتار سامانه ساخته شده
است \mdash{} و کارش را هم به‌خوبی انجام می‌دهد \mdash{} می‌تواند هم‌زمان توانایی ما در دیدن یک نقص بنیادی را
از بین ببرد. یک ارزیابی که تنها پیکربندی نهایی را می‌سنجد، نمی‌تواند بگوید نتیجه‌ی خوب از
درستی قانون کنترل آمده است یا از میراگری که رویش گذاشته شده.
